%Copyright 2021 Jean-Michel Bruel
%This program is free template, under the terms of the attached License of this 
%repo. 
%Based on https://github.com/jpeisenbarth/SRS-Tex
%Itself based on the code of Yiannis Lazarides
%http://tex.stackexchange.com/questions/42602/software-requirements-specification-with-latex
%http://tex.stackexchange.com/users/963/yiannis-lazarides
%Also based on the template of Karl E. Wiegers
%http://www.se.rit.edu/~emad/teaching/slides/srs_template_sep14.pdf
%http://karlwiegers.com
\documentclass{scrreprt}
\usepackage{method}		% Specific command from Handbook
\usepackage{minitoc}	% For table of content in chapters
\usepackage{listings}
\usepackage{underscore}
\usepackage{enumitem}
\usepackage{graphicx}
\usepackage[bookmarks=true]{hyperref}
\usepackage[utf8]{inputenc}
\usepackage[english]{babel}
\hypersetup{
	bookmarks=false,    % show bookmarks bar?
	pdftitle={Software Requirement Specification},    % title
	pdfauthor={Jean-Philippe Eisenbarth},                     % author
	pdfsubject={TeX and LaTeX},                        % subject of the document
	pdfkeywords={TeX, LaTeX, graphics, images}, % list of keywords
	colorlinks=true,       % false: boxed links; true: colored links
	linkcolor=blue,       % color of internal links
	citecolor=black,       % color of links to bibliography
	filecolor=black,        % color of file links
	urlcolor=purple,        % color of external links
	linktoc=page            % only page is linked
}%
% \def\myversion{1.24.8 }
\def\myversion{0}
\def\datecreated{Oct 5, 2026}
\date{}

% Setting up graphics path
\graphicspath{{./images/}}

% --- REQUIREMENT COUNTERS ---
\newcounter{ENV}
\newcounter{SFR}
\newcounter{SNFR}
\renewcommand{\theENV}{ENV-\reqnumber{\value{ENV}}}
\renewcommand{\theSFR}{SFR-\reqnumber{\value{SFR}}}
\renewcommand{\theSNFR}{SNFR-\reqnumber{\value{SNFR}}}

% Format numbers as 001, 002, 003, etc.
\newcommand{\reqnumber}[1]{%
	\ifnum#1<10 00%
	\else\ifnum#1<100 0%
	\fi\fi
	\number#1%
}

% --- REQUIREMENT MACRO SETUP ---
\newcommand{\requirement}[5]{%
	\par\noindent
	\begingroup
	
	% Increment the appropriate counter
	\refstepcounter{#1}%
	\edef\reqid{\csname the#1\endcsname}%
	\label{\reqid}%
	
	% Adjustable formatting
	\def\reqindent{2.3cm}
	\setlength{\parindent}{0pt}
	
	% Requirement ID and description
	\noindent
	\makebox[0pt][l]{\textbf{\reqid:}}%
	\hspace*{\reqindent}%
	\parbox[t]{\dimexpr\linewidth-\reqindent\relax}{%
		\parindent=0pt #2%
	}%
	\par\vspace{8pt}
	
	% Indented fields
	\leftskip=\reqindent
	\noindent\textbf{Rationale:} #3\par\vspace{5pt}
	\noindent\textbf{Fit Criterion:} #4\par\vspace{5pt}
	\noindent\textbf{Traceability:} #5\par
	
	\endgroup
	\vspace{0.8em}
}

% --- HOW TO USE IT ---

%\requirement{Tag/ID}
%{Description}
%{Rationale}
%{Fit Criterion}
%{Traceability tags}

%\title{%
	
	%}
\usepackage{hyperref}
\newcommand*{\nsection}[1]{
	\section*{#1}
	\addcontentsline{toc}{section}{#1}
}
%--------------------- Numbering ---------------
\usepackage{amsthm}
\usepackage{xassoccnt}
%\newtheorem{req}{Requirement}[section]        
\newtheorem{req}{Requirement}        
\theoremstyle{definition}        
\newtheorem{constraint}{Constraint}
\newtheorem{goal}{Goal}
\DeclareCoupledCountersGroup{theorems}
\DeclareCoupledCounters[name=theorems]{req,constraint,goal}
\setcounter{goal}{0}
%--------------------- Comments ---------------
\newcommand{\mynote}[3][black]{\textcolor{#1}{\fbox{\bfseries\sffamily\scriptsize{#2}}
		{\small\textsf{\emph{#3}}}}}
\newcommand{\comment}[1]{\mynote[red]{Comment:}{#1}}
%\newcommand{\comment}[1]{}
\newcommand{\notEmpty}{\comment{This chapter should not be empty!}}
%--------------------- Glossary ---------------
\newcommand{\glossaryDef}[3]{\item \hypertarget{{#2}}{\textbf{#1} - {#3}}}
\begin{document}
	
	\begin{center}
		\rule{16cm}{5pt}\vskip1cm
		\begin{bfseries}
			\Huge{SOFTWARE REQUIREMENTS\\ SPECIFICATION}\\
			\vspace{1.9cm}
			RealmForge\\
			\vspace{1.9cm}
			\Large{Team 20 \\
			Aeolus Software\\}
			\vspace{0.5cm}
		\end{bfseries}
		Andy Huynh\\
		Owen Johnson\\
		Christopher Ludwig\\
		Andreas Sideris\\
		Humna Saeed\\
		\vspace{4cm}
		Revision \myversion : \today
	\end{center}
	
	%===============================================================
	\chapter*{Revision History}
	%===============================================================
	
	\begin{center}
		\begin{tabular}{|c|c|c|c|}
			\hline
			Name        & Date          & Reason For Changes & Version\\
			\hline
			\hline
			Andy Huynh & 2026-10-05 & First Draft & 1.0 \\
			\hline
			Andy Huynh & 2026-10-13 & Second draft of G4 & 1.0 \\
			\hline
			Andreas Sideris & 2026-10-10    & P3 and P4 drafts & 1.0\\
			\hline
			Andreas Sideris & 2026-10-12    & E3 draft & 1.0\\
			\hline
			Andreas Sideris & 2026-10-13    & E2 draft & 1.0\\
			\hline
			Andreas Sideris & 2026-10-15    & AI disclosure and Reflection 
			updated & 1.0\\
			\hline
			Owen Johnson & 2026-10-15    & Reflection updated & 1.0\\
			\hline
		\end{tabular}
	\end{center}
	
	This document follows the requirements documentation structure presented in 
	the 
	\href{https://link.springer.com/content/pdf/10.1007/978-3-031-06739-6.pdf}{Handbook
	 of requirements and business analysis}, by Bertrand Meyer.
	
	%===============================================================
	\dominitoc% Initialization
	\tableofcontents
	\adjustmtc
	%===============================================================
	
	%===============================================================
	\addcontentsline{toc}{chapter}{Goals book}
	\chapter*{Goals}	
	\minitoc% Creating an actual minitoc
	%===============================================================
	
	\comment{Goals are "needs of the target organization, which the system will 
	address". While the development team is the principal user of the other 
	books, the Goals book addresses a wider audience: essentially, all 
	stakeholders.}
	
	\nsection{G.1 Context and overall objective}
	\comment{High-level view of the project: organizational context and reason 
	for building a system.} 
	
	\begin{goal}\label{goal:first}
		This is a goal example. If you need explicit (and automatic) numbering, 
		you can use the definitions in the \texttt{.tex} template.
		Is is refined by \ref{req:memo}
	\end{goal}
	
	\begin{req}\label{req:cross}
		This is a requirement example. It illustrates how numbering is 
		continuous and cross-types (if this is what you need).
	\end{req}
	
	\nsection{G.2 Current situation}
	\comment{Current state of processes to be addressed by the project and the 
	resulting system.}
	
	\begin{req}\label{req:memo}
		This is a requirement example. 
		It refines \ref{goal:first}
	\end{req}
	
	\nsection{G.3 Expected benefits}
	\comment{New processes, or improvement to existing processes, made possible 
	by the project’s results.}
	%---------------------------------
	\notEmpty{}
	%---------------------------------
	
	\nsection{G.4 Functionality overview}
	\comment{Overview of the functions (behavior) of the system. Principal 
	properties only (details are in the System book).}
	
	\input{G4FunctionalRequirements.text}
	
	\nsection{G.5 High-level usage scenarios}

	\input{SectionG5.text}
	
	\nsection{G.6 Limitations and exclusions}
	\comment{Aspects that the system need not address.}
	
	\nsection{G.7 Stakeholders and requirements sources}
	\comment{Groups of people who can affect the project or be affected by it, 
	and other places to consider for information about the project and system.}
	%---------------------------------
	\notEmpty{}
	%---------------------------------
	
	%===============================================================
	\chapter*{Environment}
	\addcontentsline{toc}{chapter}{Environment book}
	\mtcaddchapter
	\minitoc% Creating an actual minitoc
	%===============================================================
	
	\comment{The Environment book describes the application domain and external 
	context, physical or virtual (or a mix), in which the system will operate.}
	
	\nsection{E.1 Glossary}
	\comment{Clear and precise definitions of all the vocabulary specific to 
	the application domain, including technical terms, words from ordinary 
	language used in a special meaning, and acronyms.
		This chapter should not be empty!}

	\input{Glossary.text}
	
	\clearpage
	
	\nsection{E.2 Components}
	
	\input{SectionE2.text}
	
	\nsection{E.3 Constraints}
	
	\input{SectionE3.text}
	
	\nsection{E.4 Assumptions}
	\comment{Properties of the environment that may be assumed, with the goal 
	of facilitating the project and simplifying the system.}
	
	\nsection{E.5 Effects}
	\comment{Elements and properties of the environment that the system will 
	affect.}
	
	\nsection{E.6 Invariants}
	\comment{Properties of the environment that the system’s operation must 
	preserve.}
	
	%===============================================================
	\chapter*{System}
	\addcontentsline{toc}{chapter}{System book}
	\mtcaddchapter
	\minitoc% Creating an actual minitoc
	%===============================================================
	
	\comment{The System book refines the Goal one by focusing on more detailed 
	requirements about the system under development, mainly its constituents, 
	behaviors and properties.}
	
	\nsection{S.1 Components}
	\comment{Overall structure expressed by the list of major software and, if 
	applicable, hardware parts.}
	%---------------------------------
	\notEmpty{}
	%---------------------------------
	
	\nsection{S.2 Functionality}
	\comment{One section, S.2.n, for each of the components identified in S.2, 
	describing the corresponding behaviors (functional and non-functional 
	properties).}
	%---------------------------------
	\notEmpty{}
	%---------------------------------

	\nsection{S.3 Interfaces}
	\comment{How the system makes the functionality of S.2 available to the 
	rest of the world, particularly user interfaces and program interfaces 
	(APIs).}
	
	\nsection{S.4 Detailed usage scenarios}
	\comment{Examples of interaction between the environment (or human users) 
	and the system: use cases, user stories.}
	
	\nsection{S.5 Prioritization}
	\comment{Classification of the behaviors, interfaces and scenarios (S.2, 
	S.3 and S.4) by their degree of criticality.}
	
	\nsection{S.6 Verification and acceptance criteria}
	\comment{Specification of the conditions under which an implementation will 
	be deemed satisfactory.}
	
	%===============================================================
	\chapter*{Project}
	\addcontentsline{toc}{chapter}{Project book}
	\mtcaddchapter
	\minitoc% Creating an actual minitoc
	%===============================================================
	
	\comment{The Project book describes all the constraints and expectations 
	not about the system itself, but about how to develop and produce it.}
	
	\nsection{P.1 Roles and personnel}
	\comment{Main responsibilities in the project; required project staff and 
	their needed qualifications.}

	\section{P.1 Roles and personnel}

Project responsibilities are structured across the team, ensuring that every member participates in meeting, technical, and research roles while maintaining a primary focus on their assigned responsibilities.

\subsection{Andreas}
\begin{itemize}
\item \textbf{Meeting Role (Chair):} Generally runs meetings, taking lead on agenda creation and keeping the team focused to cover all action items during meeting.
\item \textbf{Technical Role (CI/CD Engineer):} Focuses on creating workflows to support CI and CD.
\item \textbf{Research Role (Qualitative Research Lead):} Conducts/designs post-test interviews, surveys and focus groups.
\item \textbf{Personal Background and Role Fit:} 
% TODO: Andreas to fill in background, skills, and role fit.
\end{itemize}

\subsection{Andy}
\begin{itemize}
\item \textbf{Meeting Role (Scribe):} Records everything discussed, updating tasks, deadlines, and next steps/action items on the GitHub Repository.
\item \textbf{Technical Role (Systems Architect):} Takes lead on architecture design and oversees development is according to the blueprint.
\item \textbf{Research Role (Ethics and Participant Coordinator):} Drafts consent forms and helps recruit study participants.
\item \textbf{Personal Background and Role Fit:} 
% TODO: Andy to fill in background, skills, and role fit.
\end{itemize}

\subsection{Chris}
\begin{itemize}
\item \textbf{Meeting Role (Progress Tracker):} Follows up on details from previous meetings and deadlines to ensure team stays on track with older tasks.
\item \textbf{Technical Role (Back-end Lead):} Leads development of the back-end of the program.
\item \textbf{Research Role (Research Design Lead):} Formulates the research questions, hypotheses, and experimental methodology for evaluating the user experience.
\item \textbf{Personal Background and Role Fit:} 
% TODO: Chris to fill in background, skills, and role fit.
\end{itemize}

\subsection{Humna}
\begin{itemize}
\item \textbf{Meeting Role (Quality/Scope Assessor):} Assesses the scope of each decision made and ensures that team stays within the desired workload.
\item \textbf{Technical Role (Front-end Lead):} Leads development of the front-end of the program.
\item \textbf{Research Role (Statistical Data Analyst):} Analyzes quantitative data to drive further decisions.
\item \textbf{Personal Background and Role Fit:} 
% TODO: Humna to fill in background, skills, and role fit.
\end{itemize}

\subsection{Owen}
\begin{itemize}
\item \textbf{Meeting Role (Risk Assessor):} Should question all assumptions made and attempt to poke holes in all decisions, should try and get team to think deeply about everything.
\item \textbf{Technical Role (Front-end Lead):} Leads development of the front-end of the program.
\item \textbf{Research Role (Data Collection Analyst):} Builds logging mechanisms into the software to capture quantitative user metrics (time-played, score, etc.).
\item \textbf{Personal Background and Role Fit:} 
% TODO: Owen to fill in background, skills, and role fit.
\end{itemize}
	
	\nsection{P.2 Imposed technical choices}
	\comment{Any a priori choices binding the project to specific tools, 
	hardware, languages or other technical parameters.}
	
	\nsection{P.3 Schedule and milestones}
	
	\input{SectionP3.text}
	
	\nsection{P.4 Tasks and deliverables}
	
	\input{SectionP4.text}
	
	\nsection{P.5 Required technology elements}
	\input{SectionP5.text}
	
	\nsection{P.6 Risks and mitigation analysis}
	\comment{Potential obstacles to meeting the schedule of P.4, and measures 
	for adapting the plan if they do arise.}
	
	\nsection{P.7 Requirements process and report}
	\comment{Initially, description of what the requirements process will be; 
	later, report on its steps.}
	
	\chapter*{Appendix}
	\addcontentsline{toc}{chapter}{Appendix}
	\mtcaddchapter
	
	\nsection{Appendix --- Generative AI Use Disclosure}
	
	\subsection*{AI Tools Used}
	
	\begin{itemize}
		\item OpenAI GPT-6
		\item Google Search AI Overview (Google Gemini)
	\end{itemize}
	
	\subsection*{How and Where Used}
	
	AI tools were used for assistance with LaTeX syntax and commands, during 
	the development of sections E2, E3, S4, P3, and P4, and when developing 
	LaTeX formatting for the document as a whole. Specifically, AI tools were 
	used for assistance in creating the LaTeX custom requirement format macro 
	used throughout the document, as well as resolving various minor formatting 
	and table of contents bugs. Additionally, these tools were also used in 
	these sections to conduct brief `sanity checks' of content in accordance 
	with Meyer's requirements document template (the template description was 
	provided). These `sanity checks' did not substitute full reviews of the 
	sections, and were more used as a brief second opinion. In these sections, 
	AI tools were not used to generate any text or wholly generate any of the 
	ideas or conclusions presented.
	
	\subsection*{Verification of AI-Generated Content}
	
	Any LaTeX syntax or code provided by AI tools was verified by testing 
	compilation in LaTeX documents, and ensuring the resulting output was 
	correctly formatted and matched the needs of the document. Conclusions from 
	brief `sanity checks' of the above sections using AI tools were 
	independently verified for correctness using course materials and 
	documentation of Meyer's requirements template, and no ideas or text were 
	newly generated by AI tools and presented in the document.
	
	\newpage{}
	
	\nsection{Appendix --- Reflection}
	
	\textbf{Reflection Questions Answered:}
	
	\begin{enumerate}
		\item What went well while writing this deliverable? 
		\item What pain points did you experience during this deliverable, and 
		how did
		you resolve them?
		\item How many of your requirements were inspired by speaking to your
		client(s) or their proxies (e.g. your peers, stakeholders, potential 
		users)?
		\item Which of the courses you have taken, or are currently taking, 
		will help
		your team to be successful with your capstone project.
		\item What knowledge and skills will the team collectively need to 
		acquire to
		successfully complete this capstone project?  Examples of possible 
		knowledge
		to acquire include domain specific knowledge from the domain of your
		application, or software engineering knowledge, mechatronics knowledge 
		or
		computer science knowledge.  Skills may be related to technology, or 
		writing,
		or presentation, or team management, etc.  You should look to identify 
		at
		least one item for each team member.
		\item For each of the knowledge areas and skills identified in the 
		previous
		question, what are at least two approaches to acquiring the knowledge or
		mastering the skill?  Of the identified approaches, which will each team
		member pursue, and why did they make this choice?
	\end{enumerate}
	
	\textbf{Andreas:}
	
	\begin{enumerate}
		\item When writing this deliverable, adopting Bertrand Meyer's 
		requirement template and completing its sections was helpful and 
		applicable to our project, since all team members were familiar with 
		the template previously, and it is provided with thorough 
		documentation. Understanding the template previously, and being able to 
		review its full documentation helped our team to understand what was 
		required to complete each section. Additionally, writing this 
		deliverable forced our team to deeply consider various details of the 
		project and requirements that the system must satisfy, which helped to 
		enable us to make important project decisions. Obviously one must 
		decide on various project details when writing a requirement 
		specification, but completing this SRS helped us to solidify and 
		exactly outline the outputs needed to be produced by the system, and 
		other constraints it must follow.
		\item One pain point I experienced during this deliverable was 
		coordinating the many sections to be written and decisions to be made 
		as a group. Ideally, all project and system planning decisions will be 
		made together, but work must be divided at some point so that each team 
		member can asynchronously complete their tasks (while reviewing and 
		discussing with the team when needed). This at times led to more 
		administrative overhead work, as various sections needed to be 
		iteratively reviewed and discussed to ensure that the document did not 
		contain major contradictions. However, this coordination work was 
		required for the success of the document, and was lessened in certain 
		scenarios by initially discussing clear expectations for various 
		sections as a group, to ensure that the team member completing a 
		specific section had the information required to do so correctly.
		\item A majority of the sections I contributed to did not involve the 
		production of formal requirements, except E.3 (Constraints), which 
		included formal environmental constraint requirements. Some of these 
		requirements (feature variation points, and user input device 
		requirements) were primarily influenced by discussions with our TA and 
		professor, using their recommendations on including feature/system 
		variability in the project, and specifically doing so by having players 
		use different input methods (controller, keyboard and mouse). 
		Additionally, the environment requirement included concerning the 
		operating system RealmForge will support was influenced by discussions 
		with peers that could be included in our user groups (those who play 
		video games and are familiar the roguelite genre, those who are not, 
		etc.). It was found (as expected) that most users tend to play PC video 
		games on machines running modern Windows operating systems, as they are 
		most commonly available and easy to configure to play popular video 
		games.
		\item Of the courses I have taken (and are currently taking), SFWRENG 
		3RA3 (Requirements) and 4HC3 (Human-Computer Interaction) are 
		particularly helpful to our team (and myself specifically) for our 
		capstone project. SFWRENG 3RA3 introduced all team members to the use 
		of Meyer's requirements template and principles of requirements 
		engineering as a whole, which helped us to complete this document, and 
		construct requirements for the system. SFWRENG 4HC3 provides our team 
		with important knowledge and techniques for effective user interface 
		and experience design, which will be particularly helpful for our 
		project, since it primarily focuses on interface design and usability 
		testing.
		\item One skill the team will collectively need to acquire (or further 
		develop) is video game development. This will include becoming more 
		familiar with our chosen game development engine and language, as well 
		as better understanding game development principles. This will be 
		especially important for myself, as I have limited experience in 
		programming for video game development. Another area of knowledge that 
		will need to be further developed for each team member concerns 
		usability testing techniques and principles, as each of us are just 
		beginning to study these concepts in SFWRENG 4HC3 (as discussed above).
		\item For improving game development skills, team members could use 
		provided game engine tutorials and documentation for the game engine 
		used, review YouTube and web article lessons on game development 
		principles, or read books that provide instruction on game development 
		principles. I will individually aim to primarily follow tutorials and 
		documentation for our chosen game engine, since the initial technical 
		skills of game development are new to me. I will also aim to improve my 
		understanding of game development principles through various tutorials 
		and lessons when possible. For developing interface design and 
		usability testing knowledge, team members could take notes and review 
		how to best apply the concepts taught in SFWRENG 4HC3 to our Capstone 
		project specifically, or consult other standards or tutorial sources 
		for interface design and usability testing. I will primarily aim to 
		review and understand how the concepts taught in 4HC3 could be applied 
		to our project, as I currently am developing my understanding of these 
		concepts well in the context of the course, and believe they could be 
		effective when applied to the development of RealmForge.
	\end{enumerate}
	
	\textbf{Owen:}
	
	\begin{enumerate}
		\item When writing this deliverable, collaborating using the Meyer template made the process much easier since
		everyone had worked with it previously in the course SFWRENG 3RA3. The division of work went smoothly at the
		beginning as well, and communication throughout the process made us get more specific with what exactly
		implementation would entail, what constraints/requirements exist, and how those will be satisfied.

		\item There were some pain points throughout the deliverable. The first one was the adjustment to using branches
		much more often (as opposed to a branch for each document, it was per section). This took a little bit of getting used to
		for me to do this without issue, but was resolved by discussing the process with group members, and getting
		comfortable with the process of branch management. Another pain point was consistent communication, as the bulk of
		the work was done during reading week, when people's schedules are at their most inconsistent. This was also made
		worse by many sections being written individually, without realizing at the time of writing that certain sections are
		similar, which created some redundant points. This was resolved by constantly discussing clear guidelines for sections to
		ensure everything is cohesive and not repetitive. This preventative measure started once the issue was noticed, and
		ensured more members knew exactly what was required of them.

		\item Basically all requirements I had a hand in writing were directly inspired by speaking to clients/stakeholders,
		although not always in the context of the project. I have had plenty of discussions about game design with friends and
		peers that apply to the project and its requirement elicitation, since a lot of friends of mine like games and are
		interested in game development like I am. Although many of these people are potential users/stakeholders of the
		project, most of these conversations either predate the project's inception or were not directly related. The exception is
		our TA and supervisor, both of whom are experienced with video games and stakeholders in the project.

		\item I can think of many courses, past and present, that will help the team's success with the capstone project. 
		\begin{itemize}
	
			\item SFWRENG 4HC3: This course is teaching about how to build effective user interfaces and experiences, which
			will help build usable knowledge for the project since UI and various design principles described in this course are
			all very important to the success of this project.

			\item SFWRENG 3RA3: As mentioned in question 1, the group all has experience with the Meyer SRS template
			introduced in this course and also gained important information about requirements from this course.

			\item SFWRENG 3S03: Software testing knowledge from this course will be integral to how we approach the
			project.

			\item SFWRENG 2XC3: This course was an introduction to many core concepts we use in this project, such as
			version control, testing, debugging, and CI/CD.
	
		\end{itemize}

		\item The main skill that the team as a whole will need to acquire and develop is video game development and its
		intricacies. Since I am somewhat knowledgeable about this, I will do what I can to teach the group as needed about the
		game engine and language we use. Another skill the team will need is setting up test suites and linting. This is at least a
		requirement for myself, as I am not very experienced with setting these things up, especially in conjunction with a game
		engine. The practice of good UI design is something being taught in 4HC3 as mentioned in Q4, and is important for the
		team to absorb from the course.

		\item Game Development skill acquisition approaches:
		\begin{itemize}
			\item Following online tutorials to understand the engine/language and game development techniques, either on
			YouTube or provided by the engine's documentation.
			\item Reading up on documentation to learn what is possible in the chosen engine and language.
			\item Discuss with peers or members of the group with more experience with game design and the chosen
			engine/language.	
		\end{itemize}

		Testing and Linting skill acquisition approaches:
		\begin{itemize}
			\item Experimenting with these tools in a neutral/sample environment.
			\item Following online tutorials to set up/use the chosen tools in the chosen game engine.
		\end{itemize}

		Good UI Design skill acquisition approaches:
		\begin{itemize}
			\item Following along with and drawing from SFWRENG 4HC3 as we are taught more (or consulting a peer in that
			class, if not taking the class yet) and participating in assignments to build skills.
			\item Experimenting with UI setup through YouTube tutorials, 4HC3 knowledge, or other means, tweaking the
			design based on what you personally like (and testing it on peers).
		\end{itemize}
	\end{enumerate}
	
	\textbf{Humna:}
	
	\begin{enumerate}
		\item
		\item 
		\item 
		\item 
		\item 
		\item 
	\end{enumerate}
	
	\textbf{Christopher:}
	
	\begin{enumerate}
		\item
		\item 
		\item 
		\item 
		\item 
		\item 
	\end{enumerate}
	
	\textbf{Andy:}
	
	\begin{enumerate}
		\item
		\item 
		\item 
		\item 
		\item 
		\item 
	\end{enumerate}
	
	
	%===============================================================
\end{document}
%===============================================================